\chapter{Conclusion and Future Scope}
\label{Chapter6}
\lhead{Chapter 6. \emph{Conclusion and Future Scope}}

In this work, theoretical validation of the theoretical design for a low cost two-tier emergency communication system was successfully designed and validated. The concept of the design could be legalized and economically applied to India by the chosen system concept of a structure employing LoRa, mesh protocol and the license-free 433 MHz band. The simulation results showed the good theoretical performance of a home-built QFH antenna with ideal performance of S11 = -21.25 dB with hemispherical radiation patterns. The implementation of the Tier 1 Civilian Node and its corresponding core network software for the solution was also successful in demonstrating its functionality with the Meshtastic platform. 

The thrust of this work is theoretical validation, the immediate priority for future work will be the transition from simulation to a physically robust prototype. This will necessitate the re-fabrication of the QFH antenna with precise assembly strategies in mind to ameliorate the fabrication issues which exist within the prototype and which caused the resonance to shift centres in the first prototype. Once this antenna design has been fabricated it can be manufactured is a practical way with its validation using a VNA for this design for its resonant frequency at 433 MHz. Once validated the potential will then be explored by way of extensive testing in real world scenarios to measure the range of the QFH equipped Base Station against the original antenna-equipped Civilian Nodes with a goal to quantify the true performance of the system.